\section{A uniform single-soft estimate and the overlapping-limit obstruction}\label{app:softremainder}

This appendix makes one part of the transfer problem quantitative. The tree soft theorem is an established result of Sen \cite{Sen17tree}; the estimate below is a consequence on compact sets of massive hard kinematics, not a replacement for that theorem or a new proof of the complete four-dimensional infrared problem. We then give an exact example explaining why the same argument does not establish the required loop remainder.

\subsection{Kinematics and the tree pole decomposition}

Use signature $(+---)$ and write a future null soft momentum as $q=\omega n$, where $n=(1,\widehat n)$, $|\widehat n|=1$. The hard momenta $p_i$ are future timelike; incoming and outgoing legs have signs $\eta_i=-1$ and $+1$, respectively. At the limiting hard configuration, $\sum_i\eta_i p_i=0$. In the convention with linear coupling $-(\kappa/2)h_{\mu\nu}T^{\mu\nu}$, the leading coefficient for one physical graviton polarisation is
\begin{equation}\label{eq:softa}
 a(n)=\frac{\kappa}{2}\mathcal M_N
       \sum_i\eta_i\frac{\epsilon_{\mu\nu}(n)p_i^\mu p_i^\nu}{p_i\cdot n}.
\end{equation}
Here $\mathcal M_N$ is the stripped hard amplitude. A smooth on-shell recoil prescription is fixed when imposing momentum conservation at nonzero $q$. Its changes at order $\omega$ belong in the remainder. Formula~\eqref{eq:softa} is unchanged by $\epsilon_{\mu\nu}\mapsto\epsilon_{\mu\nu}+n_\mu\xi_\nu+\xi_\mu n_\nu$, since the change is $\kappa\mathcal M_N\xi\cdot\sum_i\eta_i p_i=0$.

Choose a compact set of nonsingular hard configurations, away from additional hard-channel poles and thresholds, and restrict to a recoil prescription defined on a neighbourhood of it. Since each hard leg is massive,
\begin{equation}\label{eq:massivegap}
 p_i\cdot n\ge p_i^0-|\boldsymbol p_i|
 =\frac{m_i^2}{p_i^0+|\boldsymbol p_i|}>0.
\end{equation}
After decreasing the upper soft energy $E$ if necessary, the same lower bound with a smaller positive constant holds uniformly along the recoil family. A physical helicity basis can be chosen in finitely many angular charts, a finite cover of the sphere of the kind whose combinatorics is studied in \cite{BB26}; the bounds below concern its invariant norms and do not require a globally smooth polarisation vector on the sphere.

In the tree sectors covered by \cite{Sen17tree}, use the usual meromorphic tree amplitude and its external-leg pole factorisation. On the selected compact neighbourhood it has the form
\begin{equation}\label{eq:poleform}
 \mathcal M_{N+1}(\omega,n)=\frac{1}{\omega}\sum_i f_i(\omega,n)+h(\omega,n),
 \qquad \sum_i f_i(0,n)=a(n),
\end{equation}
where the $f_i$ are continuously differentiable in $\omega$ and $h$ is bounded, uniformly in the hard configurations and the angular charts. These are tree analytic inputs, not assumptions about a loop amplitude. One obtains the $f_i$ by absorbing the nonzero factor $2\eta_i p_i(\omega,n)\cdot n$ into the numerator of each external-leg propagator. Other channels are regular by the choice of neighbourhood. The amplitude may contain infinitely many string poles elsewhere; the statement uses only local regularity away from those poles, rather than a truncation of the string spectrum.

\begin{proposition}[uniform tree remainder]\label{prop:treeremainder}
Let $H=\sup|h|$, $D_i=\sup|\partial_\omega f_i|$, and $C=H+\sum_iD_i$ on the neighbourhood just specified. Then
\begin{equation}\label{eq:treeR}
 \mathcal M_{N+1}(\omega,n)=\frac{a(n)}{\omega}+r(\omega,n),
 \qquad |r(\omega,n)|\le C\quad(0<\omega\le E).
\end{equation}
\end{proposition}
\begin{proof}
Subtract the coefficient at $\omega=0$ in \eqref{eq:poleform}. The fundamental theorem of calculus gives
\[
 r(\omega,n)=h(\omega,n)
       +\sum_i\int_0^1\partial_\omega f_i(t\omega,n)\,dt.
\]
The stated bound follows by the triangle inequality. All suprema are finite by the regularity and compactness hypotheses. The lower bound \eqref{eq:massivegap} is what makes the estimate uniform in angle for this argument. It supplies no such constant for a massless hard leg.
\end{proof}

\subsection{An explicit bound after the real-emission subtraction}

Let $\Omega$ be the angular and physical-polarisation space, with finite measure $\mu(\Omega)$; for both helicities integrated over the full sphere, $\mu(\Omega)=8\pi$. Put $A=\sup_\Omega|a|$. The one-particle phase-space convention is
\[
 d\Phi_q=\frac{\omega\,d\omega\,d\mu(n)}{2(2\pi)^3}.
\]
Let a real measurement function satisfy $|F(\omega,n)|\le B$ and $|F(\omega,n)-F(0,n)|\le L\omega$. Smooth recoil Jacobians and bounded hard-bin weights may be included in $F$ when these two estimates hold.

\begin{theorem}[integrable subtracted tree emission]\label{thm:realtail}
Under Proposition~\ref{prop:treeremainder}, define
\[
 T(\omega,n)=F(\omega,n)|\mathcal M_{N+1}(\omega,n)|^2
            -F(0,n)\frac{|a(n)|^2}{\omega^2}.
\]
The signed integral $\mathcal R_\lambda(F)=\int_{\lambda<\omega<E}T\,d\Phi_q$ has a limit as $\lambda\downarrow0$, and
\begin{equation}\label{eq:realtail}
 |\mathcal R_0(F)-\mathcal R_\lambda(F)|
 \le \frac{\mu(\Omega)}{2(2\pi)^3}
 \left[(LA^2+2BAC)\lambda+\frac{BC^2}{2}\lambda^2\right].
\end{equation}
\end{theorem}
\begin{proof}
The complex amplitude identity in \eqref{eq:treeR} gives
\[
 \left|\left|\frac a\omega+r\right|^2-\frac{|a|^2}{\omega^2}\right|
 \le\frac{2|a||r|}{\omega}+|r|^2.
\]
Separate $T$ into $(F-F(0,n))|a|^2/\omega^2$ and $F$ times the difference of the two squared amplitudes. Consequently
\[
 |T(\omega,n)|\le \frac{LA^2+2BAC}{\omega}+BC^2.
\]
Multiplication by the phase-space factor $\omega$ leaves an integrable bound at zero. Integrating it over $(0,\lambda)$ proves absolute convergence and \eqref{eq:realtail}.
\end{proof}

This proves cutoff removal for the \emph{subtracted single-real-emission tree sector}. It neither assigns a finite value to its unsubtracted logarithmic term nor supplies the corresponding virtual contribution. Those are separate ingredients in an inclusive observable.

The lower bound \eqref{eq:massivegap} applies to a real on-shell null $q$, not to an arbitrary virtual loop momentum. For example $p=(m,0,0,0)$ and $k=(0,\boldsymbol k)$ give $p\cdot k=0$ even when $\boldsymbol k\ne0$. Hence there is no estimate $|p\cdot k|\ge c\|k\|$ on the complete real loop domain merely because $p$ is massive. A virtual estimate must handle the pole prescription, contour and cancellations; it cannot be obtained by copying the real-emission angular bound.

\subsection{Logarithmic remainders: a sufficient quantitative bound}

A bounded tree remainder must not be assumed unchanged at four-dimensional loop level. Logarithmic soft terms are known \cite{SahooSen19}. Their presence alone does not prevent integration after the leading subtraction. For a nonnegative integer $m$, define
\begin{align}
 h_m(\omega)&=(1+\ln(E/\omega))^m,\\
 P_m(x)&=\sum_{j=0}^m\binom mj j!x^{m-j},\\
 Q_k(x)&=\sum_{j=0}^k\binom kj\frac{j!}{2^j}x^{k-j}.
\end{align}
For $0<\lambda\le E$ and $x=1+\ln(E/\lambda)$, direct integration gives
\begin{align}\label{eq:logmoments}
 \int_0^\lambda h_m(\omega)\,d\omega&=\lambda P_m(x),\\
 \int_0^\lambda\omega h_m(\omega)^2\,d\omega
 &=\frac{\lambda^2}{2}Q_{2m}(x).
\end{align}
Indeed put $\omega=\lambda e^{-u}$, expand $(x+u)^m$, and use $\int_0^\infty e^{-cu}u^jdu=j!/c^{j+1}$, proved by $j$ integrations by parts. The exponential decay makes each boundary term vanish.

\begin{proposition}[logarithmic extension]\label{prop:logremainder}
If a specified amplitude has the decomposition \eqref{eq:treeR} with the weaker bound $|r(\omega,n)|\le Ch_m(\omega)$, the same subtraction is absolutely integrable and its tail is at most
\begin{equation}\label{eq:logtail}
 \frac{\mu(\Omega)}{2(2\pi)^3}
 \left[LA^2\lambda+2BAC\lambda P_m(x)
       +\frac{BC^2\lambda^2}{2}Q_{2m}(x)\right].
\end{equation}
If a regulator-dependent family $r_\lambda$ has this bound with the \emph{same} $C,m$ and converges pointwise to $r_0$, its subtracted integrals converge by dominated convergence. If in addition $|r_\lambda-r_0|\le\varepsilon(\lambda)h_m$ for $\omega\ge\lambda$, their difference is bounded by the tail \eqref{eq:logtail} plus
\begin{equation}\label{eq:regmodulus}
 \frac{\mu(\Omega)}{2(2\pi)^3}\varepsilon(\lambda)
 \left[2BA E P_m(1)+BC E^2Q_{2m}(1)\right].
\end{equation}
\end{proposition}
\begin{proof}
Repeat the proof of Theorem~\ref{thm:realtail}, replacing $C$ by $Ch_m(\omega)$, and then use \eqref{eq:logmoments}. These moments give a regulator-independent integrable majorant when $C,m$ are uniform. For the quantitative last assertion,
\[
 \big||a/\omega+r_\lambda|^2-|a/\omega+r_0|^2\big|
 \le\frac{2A\varepsilon(\lambda)h_m}{\omega}
       +2C\varepsilon(\lambda)h_m^2.
\]
Integrate the right side with the measurement bound $B$ and phase space, extending its domain to $(0,E)$. This proves \eqref{eq:regmodulus}.
\end{proof}

The proposition does not assert that the remainders of every string loop amplitude satisfy its uniform hypothesis. Bounds whose constant diverges with the virtual infrared regulator do not meet it.

\subsection{Why separate soft limits do not control a corner}

\begin{proposition}[an overlapping-limit obstruction]\label{prop:softcorner}
For every fixed $x>0$, the integral $\int_0^E(x+y)^{-2}\,dy$ is finite, and the analogous statement holds with $x$ and $y$ exchanged. Nevertheless,
\begin{equation}\label{eq:corner}
 \int_\lambda^E\int_\lambda^E\frac{dx\,dy}{(x+y)^2}
 =\ln\frac{(E+\lambda)^2}{4E\lambda}
 \longrightarrow+\infty.
\end{equation}
\end{proposition}
\begin{proof}
The first assertion follows from the antiderivative $-1/(x+y)$. Integrating first in $x$ gives $1/(\lambda+y)-1/(E+y)$. Integrating in $y$ yields \eqref{eq:corner}. In coordinates $r=x+y$, $t=x/(x+y)$, the Jacobian has absolute value $r$, so the density becomes $dr\,dt/r$. Resolving the corner by this change of variables exposes a logarithmic radial divergence; it does not remove it.
\end{proof}

This is an analytic counterexample to an inference from separate integrability to joint integrability. It is not asserted to be an actual string diagram or a counterexample to a gravitational Ward identity. It shows why a proof about one soft external momentum cannot replace an estimate at every overlapping degeneration.

For clarity, the useful estimate after such a resolution can also be stated exactly. Suppose one \emph{subtracted} sector of an integral is represented by variables $u_j\in(0,1]$, an angular variable $\theta$, and density
\[
 G(u,\theta)\prod_{j=1}^s\frac{du_j}{u_j}\,d\nu(\theta).
\]
If there are $\delta_j>0$ and an integrable $C(\theta)\ge0$ with
\begin{equation}\label{eq:sectorbound}
 |G(u,\theta)|\le C(\theta)\prod_j u_j^{\delta_j},
\end{equation}
then its absolute integral is at most $\|C\|_1\prod_j\delta_j^{-1}$. Removing the region where any $u_j<t$ changes the integral by at most
\begin{equation}\label{eq:sectortail}
 \|C\|_1\left(\prod_j\delta_j^{-1}\right)\sum_jt^{\delta_j}.
\end{equation}
This follows from Tonelli's theorem, $\int_0^t u^{\delta-1}du=t^\delta/\delta$, and a union bound on the removed regions. The argument proves a sufficient estimate once the sector representation and its bound are established. It does not establish \eqref{eq:sectorbound} for the string integrand or identify the physical regulator with this box cutoff.

\subsection{A useful power count for the pure-graviton soft sector}

There is also an all-graph homogeneous estimate worth separating from a full convergence theorem. Consider soft graviton propagators connecting massive eikonal lines. Suppose there are $I$ soft propagators, $V$ internal soft interaction vertices of valences $d_v\ge3$, and $E_a$ attachments to the eikonal lines. Take each interaction to have at least two soft derivatives. Exclude vacuum components and tadpoles; absorb two-point kinetic insertions into the propagator. Assume the $V$ internal momentum-conservation constraints are independent, so the number of independent soft loop momenta is $\ell=I-V$. Incidence counting gives $E_a=2I-\sum_vd_v$.

Under a uniform rescaling of all these loop momenta by $r$, the loop measures, soft propagators, eikonal denominators and vertex numerators have total degree at least
\begin{equation}\label{eq:gravdegree}
 \gamma=4\ell-2I-E_a+2V
        =\sum_v(d_v-2).
\end{equation}
Thus $V>0$ implies $\gamma\ge V>0$, whereas $V=0$ gives the marginal degree zero of independent exchanges. The equality follows by substitution of the two counting identities; the inequality uses $d_v\ge3$. If the remaining angular factors are bounded, a radial integral has density bounded by a constant times $r^{\gamma-1}dr$ and converges for $\gamma>0$.

This is consistent with the established gravitational infrared analysis \cite{ASS11}. It explains the distinguished role of independent eikonal exchanges. It is not a bound on every partial rescaling: if some momenta stay fixed, a vertex need not supply two powers of the momenta being scaled. Nor does the count bound angular pinch surfaces or other massless species. To use it in \eqref{eq:sectorbound}, every nested and overlapping scaling region must be controlled after the relevant Ward-identity cancellations. A favourable homogeneous degree for the whole graph does not prove this last assertion.

\subsection{What remains for the full string-field-theory observable}

Gravitational soft factorisation and cancellation results extend well beyond a single tree diagram \cite{Weinberg65,ASS11}. Those results must retain their scattering assumptions and the cancellations between complete sets of diagrams. They do not identify an arbitrary positive Poisson model with a complete compactified string theory containing its other massless fields.

To apply the transfer proposition of Appendix~\ref{app:infraredlimit}, one still needs a specified background, external states and infrared-safe measurement; the common real and virtual regulator; and a subtraction construction compatible with the gauge identities. At each fixed perturbative order that construction must control every simultaneous soft, collinear and degeneration region, the relevant string-state sums, and the remaining hard coefficient. A sufficient route is a uniform integrable bound after all required cancellations, such as \eqref{eq:sectorbound}, on the appropriate contours or cut phase spaces, with control of regulator removal. A rigorously controlled distributional limit could provide another route. Neither route is supplied for the complete observable here. An ordinary nonperturbative S-matrix would require further construction, beyond finiteness at each order.

The tree estimate \eqref{eq:realtail}, the logarithmic integration formulas and the corner calculation are established here. A factorisation theorem and convergent remainder for the complete four-dimensional string-field observable have not been proved. The finite algebraic certificates accompanying this appendix do not formalise the soft theorem, the phase-space integrals or the full string amplitude.
